% Figure: a single force vector resolved into components, with the
% parallelogram (polygon) construction of a resultant of two forces.
\begin{figure}[htbp]
\centering
\begin{tikzpicture}[
  force/.style={draw=engblue, -{Stealth}, very thick},
  comp/.style={draw=enggray, -{Stealth}, thin, dashed},
  res/.style={draw=engred, -{Stealth}, very thick},
  axis/.style={draw=engblack, -{Stealth}, thick},
  scale=1.0,
]
% === Left panel: one force resolved into components ===
\begin{scope}[xshift=0cm]
  \draw[axis] (0,0) -- (4.2,0) node[anchor=west, font=\small] {$x$};
  \draw[axis] (0,0) -- (0,3.4) node[anchor=south, font=\small] {$y$};
  % the force F at ~40 deg
  \draw[force] (0,0) -- (3.06,2.57) node[midway, above left, font=\small, engblue] {$\vec{F}$};
  % components
  \draw[comp] (0,0) -- (3.06,0);
  \draw[comp] (3.06,0) -- (3.06,2.57);
  \node[font=\scriptsize, below] at (1.53,0) {$F_x = F\cos\theta$};
  \node[font=\scriptsize, right] at (3.06,1.28) {$F_y = F\sin\theta$};
  % angle
  \draw[draw=enggray, thin] (0.9,0) arc (0:40:0.9);
  \node[font=\scriptsize] at (1.15,0.35) {$\theta$};
  \node[font=\small, anchor=north] at (1.8,-0.7) {(a) resolution};
\end{scope}
% === Right panel: parallelogram rule for the resultant ===
\begin{scope}[xshift=6.2cm]
  \coordinate (O) at (0,0);
  \coordinate (F1) at (3.2,0.6);
  \coordinate (F2) at (1.1,2.6);
  \coordinate (R) at (4.3,3.2);
  \draw[force] (O) -- (F1) node[midway, below, font=\small, engblue] {$\vec{F}_1$};
  \draw[force] (O) -- (F2) node[midway, left, font=\small, engblue] {$\vec{F}_2$};
  % parallelogram completion
  \draw[draw=enggray, thin, dashed] (F1) -- (R);
  \draw[draw=enggray, thin, dashed] (F2) -- (R);
  \draw[res] (O) -- (R) node[midway, above left, font=\small, engred] {$\vec{R}$};
  \node[font=\small, anchor=north] at (2.0,-0.7) {(b) parallelogram rule};
\end{scope}
\end{tikzpicture}
\caption[Force as a vector: resolution and the parallelogram rule]{A
force is a vector with magnitude, direction, and point of application.
(a) A force $\vec{F}$ at angle $\theta$ to the horizontal resolves into
Cartesian components $F_x=F\cos\theta$ and $F_y=F\sin\theta$ in the
chosen basis; the choice of basis is the engineer's and the physics is
unchanged by it. (b) Two forces sharing a point of application add by
the parallelogram rule: the resultant $\vec{R}=\vec{F}_1+\vec{F}_2$ is
the diagonal. The polygon rule (place the tail of $\vec{F}_2$ at the
head of $\vec{F}_1$) gives the same resultant and generalises to any
number of concurrent forces.}
\label{fig:vol03ch01:force-vector}
\end{figure}
